\documentclass[12pt,a4paper]{report}

% === Encoding & Language ===
\usepackage[utf8]{inputenc}
\usepackage[T2A]{fontenc}
\usepackage[bulgarian]{babel}

% === Layout ===
\usepackage[left=3cm, right=2cm, top=2.5cm, bottom=2.5cm]{geometry}
\usepackage{setspace}
\usepackage{indentfirst}
\onehalfspacing

% === Graphics & Tables ===
\usepackage{graphicx}
\usepackage{booktabs}
\usepackage{subcaption}
\usepackage{multirow}
\graphicspath{{diagrams/}}

% === Math ===
\usepackage{amsmath}
\usepackage{amssymb}

% === Bibliography ===
\usepackage{csquotes}
\usepackage[backend=biber, style=numeric, sorting=none, maxbibnames=99]{biblatex}
\addbibresource{references.bib}

% === Hyperlinks ===
\usepackage[hidelinks]{hyperref}

% === Prevent overfull hbox in long words ===
\emergencystretch=1em

% === Remove chapter numbering prefix from sections ===
\usepackage{titlesec}
\titleformat{\chapter}[display]
  {\normalfont\huge\bfseries}{\chaptertitlename\ \thechapter}{20pt}{\Huge}

\begin{document}

% ============================================================
%                        TITLE PAGE
% ============================================================
\begin{titlepage}
\begin{center}

% --- University header ---
\begin{minipage}[c]{0.15\textwidth}
    % NOTE: Place the SU logo file as images/su-logo.png
    
\includegraphics[width=\textwidth]{images/su-logo.png}
\end{minipage}%
\hfill
\begin{minipage}[c]{0.75\textwidth}
    \begin{center}
    {\large \textbf{Софийски университет ,,Св. Кл. Охридски''}}\\[4pt]
    {Факултет по математика и информатика}\\[4pt]
    {\itshape Катедра ,,Софтуерни технологии''}
    \end{center}
\end{minipage}

\vspace{1cm}
\rule{\textwidth}{0.5pt}

\vspace{1.5cm}

% --- Title ---
{\Huge \textbf{ПРЕДДИПЛОМЕН ПРОЕКТ}}

\vspace{0.8cm}

{\large на тема}

\vspace{0.5cm}

{\LARGE \textbf{,,Моделиране на семантични роли и}}\\[6pt]
{\LARGE \textbf{именувани обекти в новинарски текстове''}}

\vspace{1.5cm}

% --- Student info ---
\begin{flushleft}
\textit{Дипломант:} \textbf{Борислав Мартинов Минчев}\\[4pt]
\textit{Факултетен номер:} \textbf{5MI3400508}\\[4pt]
\textit{Специалност:} \textbf{Информатика}\\[4pt]
\textit{Магистърска програма:} \textbf{Извличане на информация и откриване на знания}
\end{flushleft}

\vspace{1.2cm}

% --- Advisor ---
\begin{flushright}
\textit{Научен ръководител:}\\[6pt]
\textbf{Проф. д-р Иван Койчев, катедра ,,Софтуерни технологии'',}\\
\textbf{ФМИ, СУ ,,Св. Климент Охридски''}

\vspace{0.8cm}

\textit{Консултант:}\\[6pt]
\textbf{докт. Димитър Димитров, катедра ,,Софтуерни технологии'',}\\
\textbf{ФМИ, СУ ,,Св. Климент Охридски''}
\end{flushright}

\vfill

% --- Footer ---
\rule{\textwidth}{0.5pt}\\[8pt]
{\large София, 2026 г.}

\end{center}
\end{titlepage}

% ============================================================
%                    TABLE OF CONTENTS
% ============================================================
\pagenumbering{arabic}
\tableofcontents
\newpage

% ============================================================
%                        CHAPTERS
% ============================================================

\chapter{Увод}

Разпространението на дезинформация и пропаганда в онлайн новинарски медии представлява едно от най-значимите предизвикателства на съвременното информационно общество. Автоматичното разпознаване на начина, по който именуваните обекти (лица, организации, държави) са представени в новинарските текстове~--- дали като протагонисти, антагонисти или невинни~--- е от ключово значение за медийния мониторинг, проверката на факти и анализа на наративни манипулации.

В контекста на споделената задача SemEval-2025 Task~10 \cite{piskorski2025semeval} е предложена формална рамка за многоезична класификация на ролите на именувани обекти в новинарски текстове. Задачата обхваща пет езика (български, английски, хинди, португалски и руски), две тематични области (войната в Украйна и климатичните промени) и три подзадачи. Настоящата работа е посветена на \textbf{подзадача~1 (Entity Framing)}, която изисква за всяко споменаване на именуван обект в дадена статия да бъде определена една основна роля (грубо ниво) и една или повече детайлни роли (фино ниво) от предварително дефинирана таксономия с 22~етикета.

Настоящият преддипломен проект поставя основите за дипломна работа по тази задача. Конкретно, целите на проекта са:
\begin{itemize}
  \item Формализиране на задачата за класификация на ролите на именувани обекти и преглед на съществуващите подходи.
  \item Задълбочен статистически анализ на многоезичния набор от данни, предоставен от организаторите на състезанието.
  \item Разработка и оценка на базов модел за класификация на грубо ниво (3~класа: Protagonist, Antagonist, Innocent), базиран на XLM-RoBERTa \cite{conneau2020xlmr} със стандартна линейна класификационна глава.
\end{itemize}

Пълната двустепенна йерархична система~--- включваща entity span pooling, семантично сходство като класификационна глава, балансирани функции на загубата, фин класификатор с меко обуславяне и адаптивно праговане~--- ще бъде разработена в рамките на дипломната работа.

Структурата на документа е следната. Глава~2 представя преглед на областта: основни понятия, съществуващи подходи и предварително обучени многоезични модели. Глава~3 описва задачата и набора от данни, включително статистически анализ и основни предизвикателства. Глава~4 представя предложения базов модел за класификация на грубо ниво: архитектура, експериментална настройка и резултати. Глава~5 обобщава изводите и очертава насоките за дипломната работа.

% ============================================================
\chapter{Преглед на областта}
% ============================================================

\section{Основни понятия}

\textbf{Рамкиране на именувани обекти} (Entity Framing) е задачата за автоматично определяне на наративната роля, която даден именуван обект изпълнява в текст~--- например дали дадено лице е представено като герой, злодей или жертва. Тази задача е тясно свързана с анализа на медийни рамки (media framing) и детекцията на пропаганда.

\textbf{Многоетикетна класификация} (Multi-label Classification) е вид класификация, при която всеки пример може едновременно да бъде свързан с няколко етикета от множество от класове. За разлика от многокласовата класификация (multi-class), където точно един етикет се присвоява на всеки пример, при многоетикетната класификация броят на присвоените етикети е променлив. В контекста на нашата задача един обект може да бъде едновременно \textit{Instigator} и \textit{Deceiver}.

\textbf{Йерархична класификация} (Hierarchical Classification) се отнася до задачи, при които етикетите са организирани в дървовидна таксономия. Предсказанията трябва да бъдат съобразени с тази йерархия~--- например, ако обект е класифициран като \textit{Guardian}, той трябва да принадлежи и към родителската категория \textit{Protagonist}.

\textbf{Трансферно обучение} (Transfer Learning) в обработката на естествен език е парадигма, при която модел първо се предварително обучава върху голям корпус от текст чрез задачи за самонаблюдение (self-supervised learning), а след това се фино настройва (fine-tuning) върху конкретна задача с по-малко данни \cite{devlin2019bert}. Този подход е в основата на съвременните езикови модели като BERT и XLM-RoBERTa.

\textbf{Представяне на именувана същност чрез специални маркери} е техника, при която в текста се вмъкват специални токени \texttt{[ENTITY]} и \texttt{[/ENTITY]} около споменаването на обекта. Моделът може след това да извлече представяне специфично за обекта чрез осредняване на скритите състояния между тези маркери (entity span pooling).


\section{Обзор на съществуващи подходи и състезания}

SemEval-2025 Task~10 \cite{piskorski2025semeval} е споделена задача, организирана в рамките на 19-ия семинар по семантична евалуация (SemEval-2025). Задачата се състои от три подзадачи: (1)~Entity Framing~--- класификация на ролята на именувани обекти; (2)~Narrative Classification~--- присвояване на наративни етикети на ниво документ; (3)~Narrative Extraction~--- генериране на текстово обяснение за доминиращия наратив. В състезанието са се регистрирали 310~екипа, от които 66~са подали официални резултати. Настоящата работа се фокусира изцяло върху подзадача~1.

Това е продължение на серията от споделени задачи за детекция на пропаганда, включително SemEval-2024 Task~4 \cite{dimitrov2024semeval} за многоезично разпознаване на техники за убеждаване в мемове.

Сред системите, постигнали най-добри резултати в подзадача~1 (Entity Framing), се открояват:

\begin{itemize}
  \item \textbf{BERTastic} \cite{mahmoud2025bertastic}~--- съчетава фино настроен XLM-RoBERTa с промптиране на GPT-4o. Постига 78,48\% точност за основната роля и 51,57\% F1 за детайлните роли. Класира се на първо място за хинди.
  \item \textbf{DUTIR} \cite{lv2025dutir}~--- използва рамка, комбинираща многоезичен превод с GLM-4-Plus, QLoRA фино настройване и ансамблова класификация. Постига първо място за английски, португалски и руски.
  \item \textbf{LTG} \cite{ronningstad2025ltg}~--- прилага евристики за избор на контекст около обекта в комбинация с XLM-RoBERTa. Постига конкурентни резултати спрямо по-големи генеративни модели.
  \item \textbf{GateNLP} \cite{singh2025gatenlp}~--- предлага йерархично тристъпково промптиране за наративна класификация. Класира се на първо място за английски.
  \item \textbf{Irapuarani} \cite{assis2025irapuarani}~--- сравнява три стратегии, комбиниращи малки и големи езикови модели. Стратегията с превод на един език показва най-добра обобщаемост.
  \item \textbf{adithjrajeev} \cite{rajeev2025adithjrajeev}~--- използва последователно обучение с резюмета, центрирани около обекта.
\end{itemize}

\begin{table}[htbp]
\centering
\caption{Сравнение на водещи системи в SemEval-2025 Task 10, подзадача~1.}
\label{tab:semeval_systems}
\begin{tabular}{llcc}
\toprule
\textbf{Система} & \textbf{Базов модел} & \textbf{Осн.\ роля (\%)} & \textbf{Fine F1 (\%)} \\
\midrule
BERTastic & XLM-R + GPT-4o & 78,48 & 51,57 \\
DUTIR & GLM-4-Plus + QLoRA & --- & --- \\
LTG & XLM-RoBERTa & --- & --- \\
GateNLP & LLM prompting & --- & --- \\
\midrule
\textbf{Нашият (baseline)} & XLM-RoBERTa-base & 68,54 & (предстои) \\
\bottomrule
\end{tabular}
\end{table}

Забележителна тенденция е, че подходите, базирани на големи езикови модели (LLM), доминират в общото класиране, но фино настроените по-малки модели (напр.\ XLM-RoBERTa) постигат конкурентни резултати при по-ниски изчислителни разходи.

\subsection{Йерархична класификация}

Йерархичната класификация на текст (Hierarchical Text Classification) е задача, при която етикетите са организирани в дървовидна таксономия с връзки ``родител--дете''. За разлика от плоската класификация, при йерархичната предсказанията на различните нива трябва да бъдат \textit{съвместими}~--- например, ако обект е класифициран с детайлна роля \textit{Guardian}, той трябва да принадлежи и към родителската категория \textit{Protagonist}.

В литературата се разграничават три основни подхода:
\begin{itemize}
  \item \textbf{Плосък подход} (flat)~--- игнорира йерархията и обучава един класификатор за всички етикети. Прост за реализация, но не гарантира съвместимост между нивата.
  \item \textbf{Локален подход} (local per-level)~--- обучава отделен класификатор за всяко ниво на йерархията. Предсказанието на горното ниво се предава като вход или ограничение на следващото. Два варианта на предаване са \textit{твърдо маскиране} (hard masking), при което недопустимите етикети се занулят, и \textit{меко обуславяне} (soft conditioning), при което вероятностите от горното ниво модулират предсказанията на долното.
  \item \textbf{Глобален подход} (global)~--- единен модел, отчитащ цялата йерархия чрез споделени представяния или рекурсивни механизми.
\end{itemize}

В настоящата работа се използва локален подход с два етапа: класификация на основната роля (грубо ниво), след което резултатът се предава на класификатора за детайлни роли (фино ниво). Този подход позволява всеки етап да бъде оптимизиран независимо, а йерархичните ограничения да бъдат наложени при прехода между нивата.

\subsection{Многоетикетна класификация}

При многоетикетната класификация (Multi-label Classification) за всеки пример може да бъде присвоено произволно подмножество от множеството от етикети. Това е по-обща задача от многокласовата класификация, при която точно един етикет се присвоява на всеки пример.

\textbf{Метрики за оценка.} Многоетикетната природа на задачата изисква специфични метрики. Основните метрики в SemEval-2025 Task~10 са:
\begin{itemize}
  \item \textbf{Exact Match Ratio (EMR)}~--- официалната метрика за класиране; измерва дела на примерите, за които предсказаният набор от етикети съвпада \textit{точно} с верния.
  \item \textbf{Micro-F1}~--- агрегира истински положителните, фалшиво положителните и фалшиво отрицателните над всички етикети и примери.
  \item \textbf{Main Role Accuracy}~--- точност на класификацията на грубо ниво (3~класа).
\end{itemize}

\textbf{Предизвикателство: дисбаланс на етикетите.} При 22~възможни етикета и средна кардиналност от~1,5, съотношението положителни/отрицателни за повечето класове е под 1:20. Стандартната бинарна кръстосана ентропия (BCE) води до модел, който минимизира загубата, като предсказва ``нищо'' за почти всички етикети. Това мотивира използването на специализирани функции на загубата: \textit{Asymmetric Loss} \cite{benbaruch2021asl} намалява приноса на лесните отрицателни примери чрез асиметрично претегляне; \textit{Focal Loss} \cite{lin2017focal} адаптивно фокусира обучението върху трудно класифицируемите примери.


\section{Многоезични предварително обучени модели}

Архитектурата Transformer \cite{vaswani2017attention}, предложена от Vaswani et al., поставя основата на съвременната обработка на естествен език. Ключовата иновация е механизмът за самовнимание (self-attention), който позволява на модела да улавя зависимости между произволни позиции в поредицата с постоянна сложност.

\textbf{BERT} (Bidirectional Encoder Representations from Transformers) \cite{devlin2019bert} адаптира кодиращата част на Transformer за предварително обучение чрез маскиране на езиков модел (MLM) и предсказване на следващо изречение (NSP). Двупосочната архитектура позволява на модела да улавя контекст и от двете страни на всеки токен, което води до значителни подобрения в широк кръг задачи.

\textbf{RoBERTa} \cite{liu2019roberta} оптимизира стратегията за предварително обучение на BERT: премахва NSP задачата, увеличава размера на партидата и обема на данните, и въвежда динамично маскиране. Авторите демонстрират, че BERT е бил значително недообучен.

\textbf{XLM-RoBERTa} \cite{conneau2020xlmr} разширява подхода на RoBERTa за 100~езика, като използва 2,5~TB данни от Common Crawl. Моделът постига значителни подобрения спрямо mBERT: +14,6\% на XNLI, +13\% F1 на MLQA и +2,4\% F1 за NER. Той покрива всичките пет езика от задачата (български, английски, хинди, португалски, руски), което го прави естествен избор за нашата многоезична класификация.

В настоящата работа използваме \texttt{xlm-roberta-base} (270M параметъра, 12~Transformer слоя, скрита размерност 768).


% ============================================================
\chapter{Описание на задачата и набора от данни}
% ============================================================

\section{Същност и цели на задачата}

Подзадача~1 (Entity Framing) от SemEval-2025 Task~10 \cite{piskorski2025semeval} се дефинира формално по следния начин: дадена е новинарска статия и списък от споменавания на именувани обекти в нея, всяко определено чрез идентификатор на документа, текстово споменаване и начална/крайна позиция. За всяко споменаване трябва да бъдат определени:
\begin{enumerate}
  \item \textbf{Основна роля} (main role)~--- точно един от трите етикета: Protagonist, Antagonist, Innocent.
  \item \textbf{Детайлни роли} (fine-grained roles)~--- една или повече от 22~етикета, организирани йерархично под трите основни роли.
\end{enumerate}

Задачата е многоезична, като данните обхващат пет езика: български (BG), английски (EN), хинди (HI), португалски (PT) и руски (RU). Тематичните области са две: войната между Русия и Украйна (URW) и климатичните промени (CC).

Официалната метрика за класиране е \textbf{Exact Match Ratio (EMR)} за детайлните роли, допълнена от micro-F1 и точност на основната роля.


\section{Таксономия на ролите}

Таксономията се състои от три основни роли и 22~детайлни роли, разпределени неравномерно (таблица~\ref{tab:taxonomy}).

\begin{table}[htbp]
\centering
\caption{Таксономия на ролите: 3 основни и 22 детайлни.}
\label{tab:taxonomy}
\small
\begin{tabular}{lll}
\toprule
\textbf{Основна роля} & \textbf{Детайлна роля} & \textbf{Описание (накратко)} \\
\midrule
\multirow{6}{*}{Protagonist (6)}
  & Guardian      & Защитник, пазител \\
  & Martyr        & Мъченик за кауза \\
  & Peacemaker    & Миротворец \\
  & Rebel         & Бунтовник, борец за промяна \\
  & Underdog      & Аутсайдер в неравна борба \\
  & Virtuous      & Нравствен, етичен лидер \\
\midrule
\multirow{12}{*}{Antagonist (12)}
  & Instigator    & Подбудител на конфликт \\
  & Conspirator   & Участник в заговор \\
  & Tyrant        & Тиранин, деспот \\
  & Foreign Adversary & Външен противник \\
  & Traitor       & Предател \\
  & Spy           & Шпионин \\
  & Saboteur      & Саботьор \\
  & Corrupt       & Корумпиран \\
  & Incompetent   & Некомпетентен \\
  & Terrorist     & Терорист \\
  & Deceiver      & Измамник \\
  & Bigot         & Фанатик, нетолерантен \\
\midrule
\multirow{4}{*}{Innocent (4)}
  & Forgotten     & Забравен от обществото \\
  & Exploited     & Експлоатиран \\
  & Victim        & Жертва \\
  & Scapegoat     & Изкупителна жертва \\
\bottomrule
\end{tabular}
\end{table}

Забелязва се силна асиметрия: категория Antagonist съдържа 12~детайлни роли~--- два пъти повече от Protagonist и три пъти повече от Innocent. Това отразява по-богатото разнообразие от негативни наративни рамки в новинарските текстове и създава допълнително предизвикателство за класификацията.


\section{Структура на данните и статистически анализ}

Данните са организирани в директории по език. За всеки език са предоставени: обучително множество с анотации (идентификатор, споменаване, позиция, основна роля, детайлни роли), валидационно множество без етикети и директория с оригиналните текстове на статиите. Всяка анотация съдържа идентификатор на документа, текстово споменаване на обект, начална и крайна позиция в текста, основна роля (една от трите) и набор от детайлни роли (подмножество от 22-те). При предобработката около споменаването на обекта се вмъкват маркери \texttt{[ENTITY]} и \texttt{[/ENTITY]}.

\subsection{Статистически анализ на разпределенията}

На фигура~\ref{fig:coarse_by_lang} е представено разпределението на трите основни роли по езици. Наблюдава се ясен дисбаланс: ролята Antagonist доминира във всички езици, следвана от Protagonist, докато Innocent е най-рядко срещаната категория. Този дисбаланс е по-изразен за хинди и португалски.

\begin{figure}[htbp]
  \centering
  \includegraphics[width=0.85\textwidth]{dataset/a2_coarse_distribution_by_language.png}
  \caption{Разпределение на основните роли по езици.}
  \label{fig:coarse_by_lang}
\end{figure}

Фигура~\ref{fig:cardinality} показва разпределението на броя детайлни етикети, присвоени на всяко споменаване на обект. Повечето обекти имат 1--2~етикета, като средната кардиналност е приблизително 1,5. Тази стойност е от значение за проектирането на системата~--- тя мотивира използването на регуларизация на кардиналността с целева стойност 1,5 в бъдещата дипломна работа.

\begin{figure}[htbp]
  \centering
  \includegraphics[width=0.85\textwidth]{dataset/a3_cardinality_distribution.png}
  \caption{Разпределение на броя детайлни етикети на обект.}
  \label{fig:cardinality}
\end{figure}

Матриците на съвместно появяване на детайлните етикети (фигури~\ref{fig:cooc_protagonist}--\ref{fig:cooc_innocent}) показват, че съвместно срещащите се двойки са предимно в рамките на една и съща основна категория. При Protagonist (фигура~\ref{fig:cooc_protagonist}) \textit{Guardian} е най-честата роля (937~примера), като се комбинира най-често с \textit{Virtuous} (116) и \textit{Rebel} (61). При Antagonist (фигура~\ref{fig:cooc_antagonist}) доминира \textit{Foreign Adversary} (946), а най-честата двойка е \textit{Instigator}~+ \textit{Foreign Adversary} (64). При Innocent (фигура~\ref{fig:cooc_innocent}) категорията \textit{Victim} е силно доминираща (1080), като често се комбинира с \textit{Exploited} (41). Междукатегорийните комбинации са редки, което подкрепя валидността на йерархичния подход.

\begin{figure}[htbp]
  \centering
  \includegraphics[width=0.75\textwidth]{dataset/a1_protagonist_cooccurrence.png}
  \caption{Съвместно появяване на детайлните етикети в категория Protagonist.}
  \label{fig:cooc_protagonist}
\end{figure}

\begin{figure}[htbp]
  \centering
  \includegraphics[width=0.95\textwidth]{dataset/a1_antagonist_cooccurrence.png}
  \caption{Съвместно появяване на детайлните етикети в категория Antagonist.}
  \label{fig:cooc_antagonist}
\end{figure}

\begin{figure}[htbp]
  \centering
  \includegraphics[width=0.75\textwidth]{dataset/a1_innocent_cooccurrence.png}
  \caption{Съвместно появяване на детайлните етикети в категория Innocent.}
  \label{fig:cooc_innocent}
\end{figure}

Фигура~\ref{fig:text_length} представя разпределението на дължината на текстовете по езици. Наблюдава се значителна вариация~--- някои документи са много дълги и надхвърлят максималната дължина от 512~токена, използвана при токенизацията, което налага стратегии за отрязване (truncation).

\begin{figure}[htbp]
  \centering
  \includegraphics[width=0.85\textwidth]{dataset/a5_text_length_distribution.png}
  \caption{Разпределение на дължината на текстовете по езици.}
  \label{fig:text_length}
\end{figure}

Фигура~\ref{fig:domain} показва разпределението по тематични области за езиците, за които тази информация е налична. Трябва да се отбележи, че идентификаторите на документите за хинди и частично за португалски не съдържат маркери за домейн, поради което тези езици не са включени в диаграмата.

\begin{figure}[htbp]
  \centering
  \includegraphics[width=0.85\textwidth]{dataset/a6_domain_breakdown.png}
  \caption{Разпределение по тематични области за езиците с налична информация.}
  \label{fig:domain}
\end{figure}


\subsection{Предизвикателства}

Анализът на данните разкрива няколко ключови предизвикателства:

\begin{enumerate}
  \item \textbf{Класов дисбаланс.} Ролята Antagonist е значително по-застъпена от останалите, а Innocent е много рядка. Това налага използването на балансирани функции на загубата~--- например Class-Balanced Cross-Entropy \cite{cui2019classbalanced}.

  \item \textbf{Разреденост на многоетикетното пространство.} При 22~възможни етикета и средна кардиналност от~1,5, съотношението положителни/отрицателни за повечето класове е под 1:20. Стандартните функции на загубата (BCE) биха довели до модел, който предсказва ``нищо'' за повечето етикети.

  \item \textbf{Многоезичност.} Петте езика имат различни типологични характеристики и писмени системи (латиница, кирилица, деванагари), което изисква многоезичен предварително обучен модел.

  \item \textbf{Йерархични ограничения.} Детайлните етикети трябва да бъдат съобразени с основната категория. Грешки на грубо ниво неизбежно водят до каскадно разпространение на грешките.

  \item \textbf{Дълги документи.} Някои статии надхвърлят лимита от 512~токена, като обектът може да е споменат в различни части на текста.

  \item \textbf{Контекстна зависимост.} Ролята на обекта зависи от контекста, в който е споменат~--- една и съща личност може да бъде представена като протагонист в една статия и като антагонист в друга.
\end{enumerate}


% ============================================================
\chapter{Базов модел за класификация на грубо ниво}
% ============================================================

В тази глава е представен базовият модел за класификация на основните роли (Protagonist, Antagonist, Innocent). Целта е да се установи начално ниво на производителност (baseline), спрямо което ще бъдат оценявани по-сложните архитектури в дипломната работа. Съзнателно са избрани най-простите решения за всеки компонент на системата~--- стандартна класификационна глава, стандартна функция на загубата и минимална предобработка~--- за да се изолира приносът на самия предварително обучен модел от приноса на архитектурните надстройки.


\section{Подготовка и предобработка на данните}

\subsection{Гранулярност на текста}

Оригиналните новинарски статии варират значително по дължина~--- от няколко изречения до над 2000~думи (виж фигура~\ref{fig:text_length}). Тъй като максималната дължина на входа при токенизация е 256~токена, подаването на цял документ би довело до загуба на контекст след отрязване. Затова текстовете са разделени на абзаци (параграфи), като за всяко споменаване на обект е избран единствено абзацът, в който се намира оригиналният span. Този подход:
\begin{itemize}
  \item ограничава дължината на входа до по-управляеми размери;
  \item запазва непосредствения контекст около обекта, който е най-информативен за определяне на ролята;
  \item избягва шум от несвързани абзаци в дълги статии.
\end{itemize}

След разделяне на абзаци оригиналните позиции на споменаванията (start/end offsets) се преизчисляват спрямо началото на съответния абзац.

\subsection{Маркиране на обекта}

Около споменаването на именувания обект в текста се вмъкват маркери \texttt{[ENTITY]} и \texttt{[/ENTITY]} чрез просто заместване на низ (string replace). За разлика от по-сложни подходи (напр.\ добавяне на специални токени към речника на токенизатора), тук маркерите се третират като обикновен текст и се разделят на подтокени от стандартния SentencePiece алгоритъм на XLM-RoBERTa.

Този избор е съзнателен~--- целта на базовия модел е да използва минимални модификации, за да се установи начално ниво. Недостатъкът е, че моделът не може да локализира точно границите на обекта, тъй като маркерите нямат специална семантика в речника на токенизатора. В дипломната работа ще бъдат добавени специални токени в речника, което ще позволи прецизно entity span pooling.

\subsection{Токенизация и кодиране}

Маркираният текст се токенизира чрез стандартния XLM-RoBERTa токенизатор с максимална дължина от 256~токена, отрязване (truncation) и допълване (padding) до фиксирана дължина. Основните роли се кодират като цели числа: Protagonist~-- 0, Antagonist~-- 1, Innocent~-- 2.

Изборът на максимална дължина от 256~токена (вместо 512) е мотивиран от два фактора: (1)~тъй като работим на ниво абзац, а не на ниво документ, повечето абзаци се побират в 256~токена; (2)~по-късата дължина позволява по-голяма партида при обучение, което стабилизира градиентите. Анализ на данните показва, че при параграфна гранулярност по-малко от 5\% от примерите надхвърлят 256~токена.

\subsection{Разделяне на данните}

Организаторите предоставят обучително и валидационно множество, но не предоставят етикети за тестовото. Затова данните са разделени така:
\begin{itemize}
  \item \textbf{Обучително множество:} 80\% от оригиналното обучително множество (4\,484~примера), получено чрез стратифицирано разделяне, запазващо пропорциите на класовете.
  \item \textbf{Валидационно множество:} останалите 20\% (1\,122~примера), използвани за мониторинг на обучението и избор на най-добрата контролна точка (checkpoint).
  \item \textbf{Тестово множество:} оригиналното валидационно множество на организаторите (604~примера)~--- тъй като етикетите му са налични, то служи за крайна оценка.
\end{itemize}

Разпределението по класове в тестовото множество е: Antagonist~--- 266~примера (44,0\%), Protagonist~--- 191~примера (31,6\%), Innocent~--- 147~примера (24,3\%).


\section{Архитектура на модела}

Базовият модел използва стандартната архитектура на HuggingFace \texttt{AutoModelForSequenceClassification} с линейна класификационна глава върху представянето на \texttt{[CLS]} токена. В следващите подсекции са описани основните компоненти на архитектурата.

\subsection{Предварително обучен модел и класификационна глава}

\paragraph{Предварително обучен модел.} Използваме \texttt{xlm-roberta-base} \cite{conneau2020xlmr}~--- многоезичен Transformer модел с 12~слоя, скрита размерност $d=768$ и общо 278~милиона параметъра. Моделът е предварително обучен чрез задача за маскиран езиков модел (MLM) върху 2,5~TB текст от Common Crawl, покриващ 100~езика, включително всичките пет езика от задачата.

\paragraph{Класификационна глава.} Стандартната глава на \texttt{XLMRobertaForSequenceClassification} се състои от два линейни слоя с активационна функция tanh между тях:
\begin{equation}
  \hat{\mathbf{y}} = \mathbf{W}_2 \, \tanh(\mathbf{W}_1 \, \mathbf{h}_{\texttt{[CLS]}} + \mathbf{b}_1) + \mathbf{b}_2
  \label{eq:classifier_head}
\end{equation}
където $\mathbf{h}_{\texttt{[CLS]}} \in \mathbb{R}^{768}$ е скритото състояние на токена \texttt{[CLS]} от последния Transformer слой, $\mathbf{W}_1 \in \mathbb{R}^{768 \times 768}$ и $\mathbf{W}_2 \in \mathbb{R}^{3 \times 768}$ са обучаеми матрици, а $\hat{\mathbf{y}} \in \mathbb{R}^3$ е изходният вектор от логити. Крайните вероятности се получават чрез softmax: $p_i = \frac{\exp(\hat{y}_i)}{\sum_j \exp(\hat{y}_j)}$.

Токенът \texttt{[CLS]} агрегира информация от целия вход чрез механизма за самовнимание, но не кодира експлицитно позицията на обекта~--- тази информация се подава имплицитно чрез текстовите маркери \texttt{[ENTITY]}/\texttt{[/ENTITY]}.

\subsection{Селективно замразяване и функция на загубата}

\paragraph{Селективно замразяване.} За да се ограничи пренастройването при малък обем данни, замразяваме всички параметри на модела с изключение на:
\begin{itemize}
  \item Слой~11 (последният Transformer блок)~--- 7,1M параметъра,
  \item Pooler слоя~--- 590K параметъра,
  \item Класификационната глава~--- 1,8K параметъра.
\end{itemize}
Това дава общо 7,68M обучаеми параметъра от 278M (2,76\%). Замразяването на слоеве 0--10 запазва общите лингвистични представяния, научени при предварителното обучение, докато слой~11 се адаптира към специфичната задача. Този подход е аналогичен на техниката за извличане на признаци с фино настройване на последния слой (\textit{feature extraction with fine-tuning of the last layer}), широко използвана при трансферно обучение \cite{devlin2019bert}.

\paragraph{Функция на загубата.} Използваме стандартна функция на загубата с кръстосана ентропия (Cross-Entropy Loss) без претегляне по класове:
\begin{equation}
  \mathcal{L}_{\text{CE}} = -\sum_{i=1}^{C} y_i \log(p_i)
  \label{eq:ce_loss}
\end{equation}
където $C=3$ е броят на класовете, $y_i$ е индикаторна функция за верния клас, а $p_i$ е предсказаната вероятност. Съзнателно е избрано да не се използва балансирана загуба (например Class-Balanced CE \cite{cui2019classbalanced} или Focal Loss \cite{lin2017focal}), за да се изолира ефектът на самия модел без допълнителни корекции.


\subsection{Хиперпараметри и обучение}

Хиперпараметрите са обобщени в таблица~\ref{tab:hyperparams}.

\begin{table}[htbp]
\centering
\caption{Хиперпараметри за обучение на базовия класификатор.}
\label{tab:hyperparams}
\begin{tabular}{ll}
\toprule
\textbf{Параметър} & \textbf{Стойност} \\
\midrule
Базов модел & \texttt{xlm-roberta-base} (278M параметъра) \\
Обучаеми параметри & 7,68M (2,76\%) \\
Замразени слоеве & Слоеве 0--10 (11 от 12) \\
Макс.\ дължина на поредицата & 256 токена \\
Гранулярност на текста & Абзац \\
Маркиране на обекта & \texttt{[ENTITY]}/\texttt{[/ENTITY]} (заместване на низ) \\
\midrule
Брой епохи & 5 \\
Размер на партидата & 8 \\
Скорост на обучение & $2 \times 10^{-5}$ \\
Тегловно затихване & 0,01 \\
Функция на загубата & Кръстосана ентропия (Cross-Entropy) \\
Класификационна глава & Линейна (pooler $\rightarrow$ dense $\rightarrow$ output) \\
\midrule
Оптимизатор & AdamW \\
Стратегия за оценка & В края на всяка епоха \\
Метрика за контролна точка & Micro-F1 \\
\bottomrule
\end{tabular}
\end{table}

Обучението е реализирано чрез HuggingFace Transformers \texttt{Trainer} API. Като най-добър модел се избира контролната точка с най-висок micro-F1 на валидационното множество. Обучението отнема приблизително 20~минути на NVIDIA GPU.


\section{Резултати и анализ}

\subsection{Обучителен процес}

Фигура~\ref{fig:train_loss} показва кривите на обучителната и валидационната загуба. Обучителната загуба намалява устойчиво от~1,1 до~0,7 в рамките на 5~епохи (2\,805~стъпки). Валидационната загуба спада от~0,84 до~0,74, като конвергира след 3-тата епоха. Разликата между двете криви показва умерено пренастройване, но без драматично разминаване~--- знак, че селективното замразяване ефективно ограничава капацитета на модела.

\begin{figure}[htbp]
  \centering
  \includegraphics[width=0.95\textwidth]{baseline/c1_baseline_training_loss.png}
  \caption{Криви на обучителната и валидационната загуба на базовия модел.}
  \label{fig:train_loss}
\end{figure}

\subsection{Общи резултати}

Резултатите на базовия модел върху тестовото множество (604~примера) са представени в таблица~\ref{tab:results}.

\begin{table}[htbp]
\centering
\caption{Общи резултати на базовия класификатор.}
\label{tab:results}
\begin{tabular}{lc}
\toprule
\textbf{Метрика} & \textbf{Стойност} \\
\midrule
Точност (Accuracy) & 68,54\% \\
Претеглена Precision & 72,40\% \\
Претеглен Recall & 68,54\% \\
Претеглен F1 & 68,22\% \\
\bottomrule
\end{tabular}
\end{table}

Моделът класифицира правилно 414 от 604~примера. Разликата между Precision (72,40\%) и Recall (68,54\%) показва, че моделът е по-склонен да пропуска примери от определени класове, отколкото да генерира фалшиви положителни резултати за тях.

\subsection{Анализ по класове}

Таблица~\ref{tab:per_class} представя детайлни метрики за всеки клас, а фигура~\ref{fig:per_class} ги визуализира.

\begin{table}[htbp]
\centering
\caption{Метрики по класове на базовия модел.}
\label{tab:per_class}
\begin{tabular}{lcccccc}
\toprule
\textbf{Клас} & \textbf{Support} & \textbf{Predicted} & \textbf{Precision} & \textbf{Recall} & \textbf{F1} \\
\midrule
Protagonist & 191 & 180 & 0,683 & 0,644 & 0,663 \\
Antagonist  & 266 & 345 & 0,638 & 0,827 & 0,720 \\
Innocent    & 147 &  79 & 0,899 & 0,483 & 0,628 \\
\bottomrule
\end{tabular}
\end{table}

\begin{figure}[htbp]
  \centering
  \includegraphics[width=0.85\textwidth]{baseline/c4_baseline_per_class_metrics.png}
  \caption{Precision, Recall и F1 по класове за базовия модел.}
  \label{fig:per_class}
\end{figure}

Наблюдават се три различни профила на грешки:

\begin{itemize}
  \item \textbf{Antagonist} показва висок Recall (82,7\%), но нисък Precision (63,8\%). Моделът е обучен в среда, в която Antagonist е мажоритарният клас (44\% от тестовите примери), и развива системно отклонение (bias) към тази категория. Предсказва Antagonist за 345~примера при реални 266~--- надвишение от 79~примера (+29,7\%).

  \item \textbf{Innocent} показва обратния модел: изключително висок Precision (89,9\%), но нисък Recall (48,3\%). Когато моделът предсказва Innocent, той е почти винаги прав, но пропуска повече от половината реални случаи. Предсказва Innocent само за 79~примера при реални 147~--- дефицит от 68~примера ($-46,3\%$).

  \item \textbf{Protagonist} е балансиран, но с умерена производителност (F1=66,3\%). Предсказва 180~примера при реални 191, което е близко до действителното разпределение.
\end{itemize}

Този модел на предсказване отразява класическо \textit{отклонение към мажоритарния клас} (majority class bias)~--- модел, обучен без балансиране на класовете, систематично предпочита по-честите категории за сметка на по-редките.

\subsection{Анализ на грешките}

Матрицата на объркване (фигура~\ref{fig:confusion}) дава пълна картина на грешките.

\begin{figure}[htbp]
  \centering
  \includegraphics[width=0.7\textwidth]{baseline/c2_baseline_confusion_matrix.png}
  \caption{Матрица на объркване за базовия класификатор (604~примера).}
  \label{fig:confusion}
\end{figure}

Общият брой грешки е 190~от 604~примера (31,5\%). Основните модели на объркване са:

\begin{enumerate}
  \item \textbf{Protagonist $\rightarrow$ Antagonist} (66~грешки, 34,7\% от грешките)~--- Това е най-честата грешка. Protagonist и Antagonist са семантично противоположни, но в новинарски текстове и двете роли се асоциират с активно агентство~--- лицата действат, взимат решения, влияят на събития. Различието е в оценъчната рамка (положителна или отрицателна), която е силно контекстно зависима и може да бъде субтилна.

  \item \textbf{Innocent $\rightarrow$ Antagonist} (59~грешки, 31,1\%)~--- Innocent обекти (жертви, забравени, експлоатирани) погрешно се класифицират като Antagonist. Тази грешка е особено проблематична, защото обръща наративната интерпретация~--- жертва се представя като злодей. Възможна причина е, че Innocent обектите често се споменават в контекст на конфликт и насилие, който моделът асоциира с Antagonist.

  \item \textbf{Antagonist $\rightarrow$ Protagonist} (40~грешки, 21,1\%)~--- обратната посока на грешка~(1), по-рядко срещана поради отклонението на модела към Antagonist.

  \item \textbf{Innocent $\rightarrow$ Protagonist} (17~грешки, 8,9\%)~--- по-рядко и по-малко критично объркване.

  \item \textbf{Грешки към/от Innocent са асиметрични:} само 8~примера (2+6) са погрешно класифицирани \textit{като} Innocent, докато 76~примера (17+59) са погрешно класифицирани \textit{от} Innocent. Моделът е прекалено консервативен при присвояване на етикет Innocent.
\end{enumerate}

\subsection{Анализ по езици}

Фигура~\ref{fig:lang_f1} и таблица~\ref{tab:per_lang} представят производителността по езици.

\begin{figure}[htbp]
  \centering
  \includegraphics[width=0.85\textwidth]{baseline/c3_baseline_per_language_f1.png}
  \caption{Претеглен F1 по езици с брой примери.}
  \label{fig:lang_f1}
\end{figure}

\begin{table}[htbp]
\centering
\caption{Детайлни резултати по езици.}
\label{tab:per_lang}
\begin{tabular}{lcccccc}
\toprule
\textbf{Език} & $n$ & \textbf{Accuracy} & \textbf{WF1} & \textbf{F1\textsuperscript{Prot}} & \textbf{F1\textsuperscript{Antag}} & \textbf{F1\textsuperscript{Inn}} \\
\midrule
PT & 116 & 81,0\% & 81,2\% & 0,812 & 0,643 & 0,893 \\
EN &  91 & 76,9\% & 72,0\% & 0,250 & 0,866 & 0,000 \\
BG &  31 & 71,0\% & 63,0\% & 0,400 & 0,816 & 0,000 \\
HI & 280 & 62,1\% & 60,7\% & 0,690 & 0,624 & 0,440 \\
RU &  86 & 62,8\% & 59,9\% & 0,531 & 0,714 & 0,182 \\
\bottomrule
\end{tabular}
\end{table}

Наблюдават се значителни разлики между езиците:

\begin{itemize}
  \item \textbf{Португалски (PT)} постига най-висок резултат (WF1=81,2\%). Причината е балансираното разпределение на класовете~--- 29~Protagonist, 28~Antagonist, 59~Innocent~--- при което моделът не страда от отклонение към мажоритарния клас. Освен това португалските текстове за климатични промени (домейн CC) съдържат ясно разграничени роли: географски обекти и екосистеми като Innocent (Pantanal, Amazónia), институции като Antagonist.

  \item \textbf{Английски (EN)} постига добра обща точност (76,9\%), но с изключително небалансирани резултати по класове: F1=0,866 за Antagonist (73~от 91~примера са Antagonist), но F1=0,000 за Innocent (нито един от 9-те Innocent примера не е разпознат правилно). Моделът практически не различава Innocent от Antagonist в английски.

  \item \textbf{Български (BG)} показва подобен модел на английски: висок Antagonist F1 (0,816), но нулев Innocent F1 (0 от 3 примера). Малкият брой BG примери (31) ограничава надеждността на оценката.

  \item \textbf{Хинди (HI)} има най-много примери (280) и показва най-балансирано представяне по класове, но с по-ниски абсолютни стойности. Хинди използва писменост деванагари, типологически различна от латиницата и кирилицата, което може да представлява допълнително предизвикателство за трансферното обучение.

  \item \textbf{Руски (RU)} е най-слабият език (WF1=59,9\%) въпреки кирилската писменост, споделена с българския. Innocent F1 е само 0,182 (1 от 9 примера), а Protagonist F1 е 0,531. Руските текстове за войната в Украйна съдържат силно зависещи от контекста роли (например Русия може да бъде Protagonist или Antagonist в зависимост от гледната точка на статията), което усложнява класификацията.
\end{itemize}

Особено забележителна е пълната неспособност на модела да разпознава Innocent примери за BG и EN~--- при тези езици 100\% от Innocent примерите са погрешно класифицирани като Antagonist. За HI процентът е 48\%, за RU~--- 89\%, а за PT~--- само 12\%. Тази систематична грешка показва, че проблемът не е само в обема на данните, а и в лексикалния и контекстуалния припокриване между Innocent и Antagonist: жертвите на конфликт се споменават в същите текстове и контексти като извършителите.

\subsection{Обобщение на резултатите}

Базовият модел постига 68,22\% претеглен F1, което е разумно начално ниво за класификация на 3~класа в многоезична среда. Анализът разкрива три основни слабости, които мотивират по-сложните подходи в дипломната работа:

\begin{enumerate}
  \item \textbf{Отклонение към мажоритарния клас} (Antagonist)~--- моделът предсказва Antagonist за 57\% от примерите при реални 44\%. Това мотивира използването на балансирани функции на загубата.
  \item \textbf{Неразличаване на Innocent}~--- при 4 от 5 езика Innocent F1 е под 0,45. Класът изисква разбиране на пасивна агентност (жертви, засегнати), което стандартният \texttt{[CLS]} токен не улавя добре. Това мотивира entity span pooling и по-прецизно представяне на обекта.
  \item \textbf{Контекстна зависимост}~--- разликата от 21,3 процентни пункта между PT (81,2\%) и RU (59,9\%) показва, че производителността зависи силно от тематичния домейн и разпределението на класовете. Това мотивира стратегии за обучение, отчитащо спецификата на домейна.
\end{enumerate}


% ============================================================
\chapter{Заключение и насоки за дипломната работа}
% ============================================================

В настоящия преддипломен проект бяха постигнати следните резултати:

\begin{enumerate}
  \item \textbf{Формализиране на задачата.} Задачата за рамкиране на именувани обекти (Entity Framing) от SemEval-2025 Task 10 беше описана формално, включително таксономията от 3~основни и 22~детайлни роли, йерархичната структура и метриките за оценка.

  \item \textbf{Анализ на данните.} Беше извършен задълбочен статистически анализ на многоезичния набор от данни (5~езика, 2~домейна), който разкри: (а)~силен класов дисбаланс с доминираща категория Antagonist; (б)~разреденост на многоетикетното пространство при средна кардиналност от~1,5; (в)~съвместно срещащите се двойки от детайлни етикети са почти изключително в рамките на една и съща основна категория, което валидира йерархичния подход.

  \item \textbf{Базов модел.} Беше разработен и оценен базов модел за класификация на грубо ниво, базиран на XLM-RoBERTa-base със стандартна линейна класификационна глава и функция на загубата с кръстосана ентропия. Моделът постига 68,22\% претеглен F1 на тестовото множество при 2,76\% обучаеми параметри.

  \item \textbf{Задълбочен анализ на грешките.} Анализът разкри систематично отклонение към мажоритарния клас Antagonist (+79~излишни предсказания), пълна неспособност за разпознаване на Innocent при 3 от 5~езика и значителна вариация по езици (от 59,9\% за руски до 81,2\% за португалски). Тези конкретни слабости насочват работата върху дипломната работа.
\end{enumerate}

\subsection*{Ограничения}

Базовият модел е съзнателно опростен, за да установи долна граница на производителността. Основните ограничения са:

\begin{itemize}
  \item \textbf{Липса на балансиране на класовете.} Стандартната функция на загубата с кръстосана ентропия не компенсира дисбаланса между Antagonist (44\%), Protagonist (32\%) и Innocent (24\%), което води до систематично предпочитане на Antagonist.
  \item \textbf{Представяне чрез \texttt{[CLS]} токен.} Токенът \texttt{[CLS]} обобщава целия вход, но не кодира експлицитно позицията на обекта. При абзаци с множество споменати лица моделът не може да разграничи кое лице е целевото.
  \item \textbf{Само грубо ниво.} Моделът класифицира само основната роля (3~класа), без детайлни роли (22~класа), което е непълно решение на задачата.
  \item \textbf{Маркиране чрез заместване на низ.} Простото заместване на текст не добавя специални токени в речника, поради което маркерите се разделят на подтокени и не могат да бъдат използвани за точно определяне на границите на обекта.
\end{itemize}

\subsection*{Насоки за дипломната работа}

Анализът на слабостите на базовия модел мотивира следните подобрения, планирани за дипломната работа:

\begin{itemize}
  \item \textbf{Entity span pooling}~--- извличане на представяне чрез осредняване на скритите състояния между маркерите \texttt{[ENTITY]} и \texttt{[/ENTITY]}, добавени като специални токени в речника. Това преодолява ограничението на \texttt{[CLS]} токена, като позволява на модела да фокусира вниманието си върху конкретния обект.

  \item \textbf{Семантично сходство като класификационна глава}~--- замяна на линейната глава с подход, базиран на косинусово подобие спрямо опорни представяния (anchor embeddings) на ролите. Този подход използва семантичното значение на имената на етикетите.

  \item \textbf{Балансирани функции на загубата}~--- Class-Balanced Cross-Entropy \cite{cui2019classbalanced} за грубото ниво и Asymmetric Loss \cite{benbaruch2021asl} за многоетикетната класификация на фино ниво. Тези функции компенсират директно мажоритарното отклонение, наблюдавано в базовия модел.

  \item \textbf{Фин класификатор} за многоетикетна класификация с 22~детайлни роли, представляващ втория етап от йерархичната система.

  \item \textbf{Меко йерархично обуславяне} (soft conditioning)~--- обучаема матрица на сродство (affinity matrix), която преобразува вероятностите от грубия класификатор в меко априорно разпределение за детайлните роли. За разлика от твърдото маскиране, този подход позволява несигурността от грубото ниво да се разпространява в системата.

  \item \textbf{Регуларизация на кардиналността}~--- квадратична наказателна функция, ограничаваща очаквания брой предсказани етикети до наблюдаваната средна стойност от~1,5 (фигура~\ref{fig:cardinality}).

  \item \textbf{Адаптивно праговане}~--- динамична стратегия за определяне на прага за многоетикетната класификация, базирана на съотношението на разликата (gap ratio) спрямо най-високата предсказана вероятност.
\end{itemize}

Настоящият преддипломен проект изгражда фундамента~--- формалната рамка, разбирането на данните и базовото ниво на производителност~--- върху който ще бъде изградена пълната йерархична класификационна система в дипломната работа.


% ============================================================
\chapter*{Използвана литература}
\addcontentsline{toc}{chapter}{Използвана литература}
\printbibliography[heading=none]

\end{document}
